\Askhsh{18235}{4}
\begin{ylh}{1}
\item 2.1 Μονοτονία-Ακρότατα-Συμμετρίες Συνάρτησης
\item 2.2 Κατακόρυφη-Οριζόντια Μετατόπιση Καμπύλης
\item 5.1 Εκθετική συνάρτηση
\item 5.2 Λογάριθμοι
\end{ylh}
Στο παρακάτω σχήμα δίνεται η γραφική παράσταση $C_f$ της συνάρτησης $f(x) = |e^x - 1|, x \in \mathbb{R}$.
\begin{center}
\begin{tikzpicture}[scale=1.2, >=latex]
    \draw[very thin, color=gray!30] (-4.5,-0.5) grid (2.5,4.5);
    \draw[->] (-4.5,0) -- (3,0) node[right] {$x$};
    \draw[->] (0,-0.5) -- (0,5) node[above] {$y$};
    \foreach \x in {-4,-3,-2,-1,1,2}
        \draw (\x,2pt) -- (\x,-2pt) node[below, font=\small] {\x};
    \foreach \y in {1,2,3,4}
        \draw (2pt,\y) -- (-2pt,\y) node[left, font=\small] {\y};
    \node[below left, font=\small] at (0,0) {$0$};
    % f(x) = |exp(x) - 1|
    \draw[very thick, color=maincolor, samples=100, smooth, domain=-4.5:1.7] plot (\x, {abs(exp(\x) - 1)}) node[above right] {$C_f$};
    % Points
    \filldraw (0,0) circle (1.5pt);
\end{tikzpicture}
\end{center}
\begin{alist}
\item Να γράψετε τον τύπο της χωρίς το σύμβολο της απόλυτης τιμής και να περιγράψετε πως αυτή μπορεί να προκύψει από τη γνωστή γραφική παράσταση της $g(x) = e^x, x \in \mathbb{R}$.
\monades{7}
\item Με τη βοήθεια της γραφικής παράστασης, ή με οποιονδήποτε άλλο τρόπο, να συμπεράνετε τη μονοτονία και την ελάχιστη τιμή της f .
\monades{6}
\item Να λύσετε την εξίσωση $f(x) = \frac{1}{2}$.
\monades{5}
\item Να βρείτε, για τις διάφορες τιμές του α, το πλήθος των κοινών σημείων της γραφικής της παράστασης $C_f$ με την ευθεία $y = \alpha$.
\monades{7}
\end{alist}

